\section{Related Work}

Our work draws on three bodies of literature that, taken together, reveal a gap: mixed reality is increasingly encountered in transient contexts, optical wearing has always been temporal but remains untheorized in wearable computing, and the cyclic nature of temporal interaction has been overlooked by transition frameworks. We review each strand before articulating how temporal wearing unifies them.

\subsection{Transient Mixed Reality in Public Contexts}

Mixed reality, as encountered by the public, is overwhelmingly transient. Museum and gallery deployments, exhibition demonstrations, pop-up VR experiences, and educational AR encounters all share a common temporal structure: visitors engage for seconds to minutes, not hours \cite{vomlehn2001exhibiting, li2024museum}. Ethnographic studies of museum conduct reveal that interactions with exhibits are brief, episodic, and socially organized---visitors negotiate access to objects, point things out to companions, and move on after dwelling for an average of 17--30 seconds \cite{vomlehn2001exhibiting}. Even dedicated exhibition VR installations, where visitors queue for a curated experience, typically offer encounters of one to five minutes before the headset is removed and the next visitor takes a turn.

Yet the devices deployed in these contexts are designed for sustained wearing. Head-mounted displays---whether tethered VR headsets at exhibition stations or AR glasses in museum pilots---are optimized for session-length use: their weight distribution, strap systems, and thermal management assume minutes-to-hours of continuous wear \cite{starner2001challenges}. This creates a practical mismatch. Visitors must navigate complex donning procedures (adjusting interpupillary distance, tightening straps, positioning displays) for encounters that may last only seconds. The ratio of setup time to experience time can approach or exceed 1:1, undermining the spontaneity that makes transient MR encounters compelling.

The mismatch extends into social dynamics. A substantial body of research documents the social costs of wearing head-mounted devices in public. Koelle et al.'s comprehensive survey of social acceptability methods in HCI identifies device visibility, perceived purpose, and wearing context as key determinants of bystander comfort \cite{koelle2020social}. Studies of data glasses adoption show that even extended familiarization fails to fully overcome negative attitudes toward always-worn face-covering devices \cite{koelle2017acceptability}. Profita et al. demonstrate that body location significantly affects the perceived social acceptability of on-body technology, with head and face locations provoking the strongest reactions \cite{profita2013wrist}. Even when HMDs serve clearly beneficial purposes---such as assistive technology for people with visual impairments---social tensions persist, with wearers reporting feelings of conspicuousness and bystanders expressing discomfort \cite{profita2016at, goodman2019social}.

Critically, these social costs are modulated by duration. A two-second peek through a museum viewfinder and a two-hour VR session at a caf\'{e} occupy fundamentally different positions on the social acceptability spectrum. Koelle et al.'s field evaluation of wearable camera acceptability found that bystander comfort depends heavily on perceived recording duration \cite{koelle2019camera}. Kelly and Gilbert's framework identifies the wearer, the device, and its use as determinants of acceptability, but does not formalize duration as a distinct variable \cite{kelly2018wearer}. The social acceptability literature has primarily examined continuous wearing scenarios, leaving the acceptability profile of brief, purposeful wearing largely unexplored.

This creates a paradox: the contexts where MR is most commonly encountered by the public---museums, exhibitions, demonstrations---are precisely the contexts where sustained-wear device assumptions are least appropriate. Transient encounters are forced into devices designed for sessions, producing friction in both the physical and social senses of the word.

\subsection{Temporal Wearing and Wearability}

The wearable computing field has developed sophisticated frameworks for understanding \textit{where} and \textit{how} devices should be worn, but has paid remarkably little attention to \textit{for how long}. Gemperle et al.'s foundational work on wearability at ISWC '98 proposed thirteen design guidelines organized around the spatial relationship between body and device: placement, form language, movement accommodation, sizing, attachment, weight distribution, and aesthetics \cite{gemperle1998wearability}. These guidelines have shaped two decades of wearable design, yet they implicitly assume extended wearing. The thirteenth guideline, ``long-term use,'' addresses durability and hygiene for prolonged wear; no corresponding guideline addresses the design challenges of \textit{brief} or \textit{transitional} wearing.

Zeagler's body maps for wearable placement extended Gemperle's framework with functional, technical, and social considerations for each body location \cite{zeagler2017where}. Social factors touch on duration indirectly---a wrist-worn device is more socially acceptable partly because wrist interactions are brief and familiar---but duration itself is not formalized as a design dimension. Starner's influential characterization of wearable computing challenges distinguishes always-available from task-specific devices \cite{starner2001challenges}, coming closest to a temporal distinction, but task-specific still implies session-length engagement of minutes to hours. Knight et al.'s comfort assessment framework includes ``ease of donning and doffing'' as a usability factor \cite{knight2006comfort}, but treats it as friction to be minimized---an overhead cost---rather than a design opportunity.

The closest precedent to temporal wearing as a design concept comes from Khurana et al.'s work on the detachable smartwatch \cite{khurana2019detachable}. Their IMWUT paper explores a smartwatch that can transition between wrist-worn and handheld modes, demonstrating that the same device can operate in different wearing modalities. This implicitly acknowledges that wearing relationships are not fixed but dynamic. However, Khurana et al. focus on device versatility---a single device serving multiple form factors---rather than on wearing duration as a design dimension. The temporal dimension of when and for how long the device is in each mode remains unexplored.

What makes this gap particularly striking is that temporal wearing has deep roots in optical culture. Long before wearable computing, humans designed and used optical devices intended for seconds-long encounters. The lorgnette, popular from the 18th through early 20th centuries, consisted of a pair of lenses mounted on a handle. Users raised it to their eyes for a focused look---at a companion across a ballroom, at a painting in a gallery, at an actor on stage---and lowered it. The gesture of raising and lowering was socially codified, even ritualized. Opera glasses served an analogous function in theatrical contexts. The monocle, the hand-held magnifying glass, the stereoscope and its parlor-room cousin the View-Master, coin-operated binoculars at scenic overlooks, and the camera viewfinder: all are temporal wearing devices. They were designed not for sustained use but for a rhythm of engagement---raise, look, lower, move, raise again.

These historical form factors embody design principles that contemporary wearable computing has largely forgotten. They optimize for instant-on activation (no boot sequence, no calibration), rapid body coupling (raise to eyes in under one second), social legibility (the gesture of looking through a device is universally understood), and graceful doffing (lowering the device is as natural as raising it). They treat the don/doff cycle not as a cost but as the fundamental unit of interaction. A lorgnette user does not ``set up'' and ``tear down'' a viewing session; each raise-and-lower \textit{is} the session.

Park et al.'s work on comfort and user acceptance establishes that acceptable weight, volume, and contact area thresholds vary by body location \cite{park2019comfort}. Viewed through a temporal lens, these thresholds also vary by \textit{duration}: users tolerate substantially more weight and bulk for a three-second peek than for a three-hour session. This suggests that temporal wearing relaxes certain comfort constraints while introducing others---particularly the need for instant activation and information density within the first second of wear.

The gap, then, is this: wearability research has extensively theorized the spatial dimension of wearing (where on the body, how attached, what form factor) but has not theorized the temporal dimension (for how long, how often, and crucially, whether the wearing duration itself constitutes the interaction). Temporal wearing fills this gap by proposing duration as a primary design variable and identifying a category---seconds-long wearing where don/doff is the interaction---that existing taxonomies do not accommodate.

\subsection{Transitional Interaction as a Cyclic Process}

Temporal wearing is fundamentally about transitions: transitions between not-wearing and wearing, between unaugmented and augmented perception, between individual and shared experience. Several research traditions have examined transitional phenomena in HCI, but none has captured the cyclic, repeating character of temporal wearing.

Weiser and Brown's vision of calm technology proposed that the most effective technologies ``engage both the center and the periphery of our attention, and in fact move back and forth between the two'' \cite{weiser1996calm}. This center-periphery model operates at the information level: a stock ticker in peripheral vision, a notification that briefly demands focal attention. Temporal wearing introduces a \textit{device-level} analog to this information-level transition. The device itself moves between absence (resting on a surface, hanging from a neck strap) and presence (raised to the eyes, held against the body). It is calm in Weiser and Brown's sense---it demands attention only briefly and purposefully, then recedes---but the transition is more dramatic than they envisioned: not a shift of attention within a persistent information environment, but a physical engagement and disengagement with a perceptual device.

Bakker et al.'s work on peripheral interaction formalizes this attention dimension \cite{bakker2015peripheral}. They characterize peripheral interaction as occurring in the ``background or periphery of attention,'' with low attention demand and seamless integration with primary activities. Temporal wearing creates a complementary pattern: a \textit{deliberate, bounded transition} from peripheral to focal attention and back. Unlike peripheral interaction, which aims to keep information in the periphery, temporal wearing momentarily foregrounds the device and its augmented content, then returns to a peripheral or absent state. The result is a punctual attention shift rather than a sustained one---what we might call a \textit{perceptual punctuation mark} in the flow of experience.

The concept of micro-mobility, introduced by Luff and Heath in the context of collaborative work, describes the fine-grained, moment-by-moment repositioning of artifacts---documents, objects, tools---within a shared workspace \cite{luff1998mobility}. People tilt, rotate, push, and pull objects to facilitate shared viewing and collaborative action. Heath et al. extend this to technology-mediated settings, examining how devices feature in the contingent, situated conduct of everyday interaction \cite{heath2000technology}. Temporal wearing extends micro-mobility from documents and artifacts to \textit{devices and the body}: the act of raising a neck-hung viewer to the eyes, tilting it for a companion to see, passing a handheld device to another person, and setting it back on a surface are all micro-mobility actions with a temporal wearing device. The device's position relative to the body is not fixed but continuously negotiated through small physical adjustments.

Recent work on transitional and hybrid interfaces has examined how users move between devices, modalities, and interaction states. Hubenschmid et al.'s survey of hybrid user interfaces in MR catalogs patterns of complementary use across 2D devices and MR environments, identifying sequential, parallel, and complementary transition types \cite{hubenschmid2025hybrid}. Ens et al.'s work on candid interaction addresses the visibility and hiddenness of wearable computing activities, proposing techniques for managing the social legibility of device use \cite{ens2015candid}. Dhaka et al. examine AR user interface transitions at ISMAR, focusing on how AR interfaces transition between visual states \cite{dhaka2024transitions}. These contributions address important aspects of transitional interaction, but they share a common assumption: transitions are \textit{linear}---from state A to state B, from device 1 to device 2, from hidden to visible. They do not account for transitions that are \textit{cyclic}---from not-wearing to wearing to not-wearing to wearing again---where the cycle itself is the fundamental unit of the experience.

Suchman's foundational work on situated action provides a theoretical grounding for why temporal wearing interactions resist linear modeling \cite{suchman2007human}. In Suchman's framework, action is not the execution of pre-formed plans but emerges from the contingent interaction between a person and their context. Temporal wearing is paradigmatically situated: the decision to pick up a device is occasioned by noticing an interesting exhibit, being prompted by a companion, or seeing another visitor's reaction. Each wearing episode is a fresh response to contextual circumstances, not a step in a predetermined sequence. This situatedness distinguishes temporal wearing from habitual wearing (permanent wear, where the device is donned once and forgotten) and planned wearing (session wear, where the user commits to an extended experience in advance).

\subsubsection{The Gap: Cyclic Transition in Temporal Wearing.}
The transition literature reveals a consistent pattern: transitions are modeled as movements between two states---center and periphery \cite{weiser1996calm}, focal and peripheral attention \cite{bakker2015peripheral}, one device and another \cite{hubenschmid2025hybrid}, visible and hidden \cite{ens2015candid}. What is missing is a model of transition as a \textit{cycle}: a repeating rhythm of engagement and disengagement where each cycle is complete in itself yet invites the next. Temporal wearing introduces precisely this cyclic structure. The interaction is not attract $\rightarrow$ don $\rightarrow$ use $\rightarrow$ doff (a linear sequence) but attract $\rightarrow$ don $\rightarrow$ glimpse $\rightarrow$ doff $\rightarrow$ rest $\rightarrow$ return (a cycle that may repeat multiple times). The ``rest'' and ``return'' phases---where the user sets the device down but remains nearby, processes what they saw, and decides whether to engage again---are as integral to the interaction as the glimpse itself.

\subsection{Summary: Three Gaps, One Concept}

Three observations emerge from this review, each supported by extensive prior work yet never connected:

\begin{enumerate}
    \item \textbf{MR is increasingly transient.} The most common public encounters with mixed reality---in museums, exhibitions, demonstrations, and educational settings---are measured in seconds, yet the devices serving these encounters are designed for sustained sessions \cite{vomlehn2001exhibiting, koelle2020social, starner2001challenges}.
    \item \textbf{Optical wearing has always been temporal.} Centuries of optical device design---lorgnettes, opera glasses, stereoscopes, viewfinders---embody a seconds-long wearing pattern, but wearable computing's spatial frameworks have not theorized this temporal dimension \cite{gemperle1998wearability, zeagler2017where}.
    \item \textbf{The interaction is cyclic, not linear.} Users engage with temporal wearing devices through repeating cycles of donning and doffing, but existing transition frameworks model only linear state changes \cite{weiser1996calm, bakker2015peripheral, hubenschmid2025hybrid}.
\end{enumerate}

\textit{Temporal wearing} unifies these three observations into a single concept. It names a widespread but untheorized practice, fills a gap in wearability taxonomies, and introduces a cyclic model of transitional interaction. In the following section, we formalize this concept, propose a wearing duration taxonomy, and define a design space for temporal wearing devices.
